% Part: second-order-logic
% Chapter: sol-and-sets
% Section: cardinalities

\documentclass[../../../include/open-logic-section]{subfiles}

\begin{document}

\olfileid{sol}{set}{crd}

\olsection{সেটের অঙ্কবাচকতা}

\begin{explain}
সংজ্ঞাক্ষেত্র যে সসীম বা অসীম, !!{enumerable} বা !!{nonenumerable}, তা যেমন
প্রকাশ করা যায়, তেমনই~$\Domain{M}$-এর কোনো উপসেট সসীম বা অসীম,
!!{enumerable} বা !!{nonenumerable} হওয়ার ধর্মও সংজ্ঞায়িত করা যায়।
\end{explain}

\begin{prop}
নিচের সূত্রটি $\fn{Inf}(X) \ident$
\begin{multline*}
\lexists[u][(\lforall[x][(X(x) \lif X(u(x)))] \land {}\\
  \lforall[x][\lforall[y][((X(x) \land X(y) \land
      \eq[u(x)][u(y)]) \lif \eq[x][y])]] \land {}\\
  \lexists[y][(X(y) \land \lforall[x][(X(x)
      \lif \eq/[y][u(x)])]])]
\end{multline*}
কোনো চলরাশি-আরোপ~$s$-এর সাপেক্ষে পরিতৃপ্ত হয় যদি এবং কেবল যদি $s(X)$
অসীম।
% BN-SRC-261 TARGET-CORRECTION-BEGIN
\par\smallskip\noindent\textnormal{\small\textbf{উৎস-সংশোধন (BN-SRC-261):} হিমায়িত সূত্রে~$u$ যে~$X$-কে~$X$-এ পাঠায় তা বলা নেই, আর একৈকতা অপ্রয়োজনীয়ভাবে গোটা সংজ্ঞাক্ষেত্রে চাপানো হয়েছে। ফলে একটি সসীম~$X$-কেও বাইরে পাঠিয়ে সূত্রটি সত্য করা যায়। বাংলা সূত্রে~$X(u(x))$-এর বদ্ধতা যোগ করা হয়েছে এবং একৈকতা~$X$-এর উপাদানযুগলে সীমাবদ্ধ করা হয়েছে, যাতে ঠিক~$X$-এর উপর একটি একৈক অসমাপতিত স্ব-অপেক্ষক দাবি করা হয়।} % BN-SRC-261 TARGET-CORRECTION-END
\end{prop}

\begin{prop}
নিচের সূত্রটি $\fn{Count}(X) \ident $
\begin{multline*}
\lnot\lexists[x][X(x)] \lor {}\\
\lexists[z][\lexists[u][(X(z) \land
    \lforall[x][(X(x) \lif X(u(x)))] \land {}\\
    \lforall[Y][(((Y \subseteq X \land Y(z)) \land
      \lforall[x][(Y(x) \lif Y(u(x)))]) \lif X = Y)])]]
\end{multline*}
কোনো চলরাশি-আরোপ~$s$-এর সাপেক্ষে পরিতৃপ্ত হয় যদি এবং কেবল যদি $s(X)$
!!{enumerable}।
% BN-SRC-262 TARGET-CORRECTION-BEGIN
\par\smallskip\noindent\textnormal{\small\textbf{উৎস-সংশোধন (BN-SRC-262):} হিমায়িত সূত্রে প্রতিটি~$u$-বদ্ধ~$Y$-কে~$X$-এর সমান বলা হয়েছে; কিন্তু~$Y$ হিসেবে গোটা সংজ্ঞাক্ষেত্র নিলেই তা সবসময় পূর্বগটি মেটায়, ফলে সূত্রটি কোনো প্রকৃত উপসেটকে তালিকায়নযোগ্য বলতে পারে না। সেটি শূন্য সেটকেও বাদ দেয়, যদিও হিমায়িত গ্রন্থের সংজ্ঞায় শূন্য সেট তালিকায়নযোগ্য। বাংলা সূত্রে শূন্য ক্ষেত্রটি পৃথক করা হয়েছে,~$u$-কে~$X$-এ বদ্ধ রাখা হয়েছে এবং ন্যূনতমতার পরিমাণায়ন~$Y \subseteq X$-এ সীমাবদ্ধ করা হয়েছে।} % BN-SRC-262 TARGET-CORRECTION-END
\end{prop}

কান্টরের উপপাদ্য থেকে আমরা জানি যে !!{nonenumerable} সেট আছে, এবং বস্তুত
অসীম আকারের অসীমসংখ্যক পৃথক স্তর আছে। সেটতত্ত্বে সেটের আকারের একটি
পূর্ণাঙ্গ পাটিগণিত গড়ে তোলা হয় এবং সেটগুলিকে অসীম অঙ্কবাচক সংখ্যা দেওয়া
হয়। স্বাভাবিক সংখ্যাগুলি সসীম সেটের আকার-মাপক অঙ্কবাচক সংখ্যা হিসেবে কাজ
করে। !!{denumerable} সেটগুলির অঙ্কবাচকতা প্রথম অসীম অঙ্কবাচকতা; একে
$\aleph_0$ (``আলেফ-নট'' বা ``আলেফ-শূন্য'') বলা হয়। পরের অসীম আকারটি
$\aleph_1$। কোনো সেট অসীম তালিকায়নযোগ্য না-হয়ে যে ক্ষুদ্রতম আকারের হতে
পারে, এটি সেই আকার (অর্থাৎ, তার আকার~$\aleph_0$ নয়)। ``$X$-এর আকার
$\aleph_0$''-কে আমরা $\fn{Aleph}_0(X) \liff \fn{Inf}(X) \land
\fn{Count}(X)$ হিসেবে সংজ্ঞায়িত করতে পারি। $X$-এর আকার~$\aleph_1$ যদি
এবং কেবল যদি সেটি অসীম, তার সব উপসেট সসীম অথবা তাদের আকার~$\aleph_0$,
কিন্তু তার নিজের আকার~$\aleph_0$ নয়। তাই নিচের সূত্রে এটি প্রকাশ করা যায়:
$\fn{Aleph_1}(X) \ident \fn{Inf}(X) \land
\lforall[Y][(Y \subseteq X \lif
  (\lnot\fn{Inf}(Y) \lor \fn{Aleph}_0(Y)))] \land \lnot
\fn{Aleph}_0(X)$।
% BN-SRC-263 TARGET-CORRECTION-BEGIN
\par\smallskip\noindent\textnormal{\small\textbf{উৎস-সংশোধন (BN-SRC-263):} হিমায়িত সংজ্ঞায়~$X$ অসীম হওয়ার দফাটি নেই। ফলে প্রতিটি সসীম~$X$-এর সব উপসেট সসীম এবং~$X$-এর আকার~$\aleph_0$ নয় বলে তাকেও ভুলভাবে~$\aleph_1$ আকারের বলা হয়। বাংলা পাঠে~$\fn{Inf}(X)$ যোগ করে ঠিক পরের অসীম অঙ্কবাচকতার অভিপ্রেত শর্ত পুনরুদ্ধার করা হয়েছে।} % BN-SRC-263 TARGET-CORRECTION-END
$\aleph_2$ আকারের হওয়া একইভাবে সংজ্ঞায়িত হয়, ইত্যাদি।

বিশেষ আগ্রহের একটি আকার আছে, যাকে সতত সমষ্টির অঙ্কবাচকতা বলা হয়। এটি
$\Pow{\Nat}$-এর আকার, অথবা তুল্যভাবে~$\Real$-এর আকার। কোনো সেট যে সতত
সমষ্টির আকারের, তা-ও দ্বিতীয়-ক্রমের যুক্তিবিদ্যায় প্রকাশ করা যায়, তবে
তার জন্য আর-একটু কাজ দরকার।

\end{document}
